\section{Introduction} \label{sec:introduction}

RDF is the data model underlying the Semantic Web and Knowledge Graphs.
 Describing and validating RDF data has received growing attention
 over the last decade, with two main technologies proposed:
   SHACL~\cite{SHACL} and ShEx~\cite{PGS14}.
 SHACL (Shapes Constraint Language) is a W3C recommendation whose definition
  was based on an informal semantics in natural language that left several features like recursion undefined.
 On the other hand, ShEx followed a community-driven approach~\cite{shex_standard} 
  whose semantic specification was based  
  on the formal semantics presented in several papers:
  \cite{SBG15} introduced the complexity of the basic validation algorithm, 
  \cite{BGP17} presented the semantics of the ShEx schema language 
   with a formal semantics that supported recursion and negation 
   based on stratified semantics, 
   and~\cite{ShExInheritance} 
   extended the semantics to support inheritance in ShEx schemas.

 Lean 4~\cite{Moura2021} has been introduced as an automatic theorem prover based on dependent type theory
   that can be used to verify the correctness of mathematical proofs.
 In this paper we present a Lean 4 representation of the ShEx
   language based on the formal semantics presented in~\cite{BGP17} and~\cite{ShExInheritance}.
 Our goal is to verify the correctness of the lemmas, propositions and algorithms presented in those papers.
 This exercise is a first step toward evaluating the suitability of automatic theorem provers
 for verifying the semantics of RDF validation languages.
 The proofs and definitions in Lean follow those given in the original papers.
 
 The main contribution of the paper is a Lean 4 representation of the ShEx language as well as a verification of the 
  correctness of the validation algorithms.
 
